\begin{problem}[Conditional probability from a class survey]
In a class, $18$ students study Mathematics Extension~1, $14$ study Chemistry, and $8$ study both subjects. One student is selected at random from the class, knowing that the student studies Chemistry. Find the probability that the student also studies Mathematics Extension~1.
\end{problem}

\begin{solution}
Let
\[
M=\{\text{studies Extension~1}\},
\qquad
C=\{\text{studies Chemistry}\}.
\]

We are asked to find
\[
P(M\mid C).
\]

Using conditional probability,
\[
P(M\mid C)=\frac{P(M\cap C)}{P(C)}.
\]
Because every student is equally likely to be chosen, we may use counts:
\[
P(M\mid C)=\frac{8}{14}=\frac47.
\]
\end{solution}

\begin{takeaways}
\item ``Given that'' signals conditional probability.
\item Once the condition is imposed, the reduced sample space becomes the relevant denominator.
\item Counts often make conditional probability faster than decimal probabilities.
\end{takeaways}

\begin{problem}[Checking independence]
A fair coin is tossed twice. Let
\[
A=\{\text{exactly one head}\},
\qquad
B=\{\text{the first toss is a head}\}.
\]
Determine whether $A$ and $B$ are independent.
\end{problem}

\begin{solution}
The sample space is
\[
\{HH,HT,TH,TT\}.
\]

Event $A$ is
\[
A=\{HT,TH\},
\]
so
\[
P(A)=\frac24=\frac12.
\]

Event $B$ is
\[
B=\{HH,HT\},
\]
so
\[
P(B)=\frac24=\frac12.
\]

The intersection is
\[
A\cap B=\{HT\},
\]
thus
\[
P(A\cap B)=\frac14.
\]

Now
\[
P(A)P(B)=\frac12\cdot\frac12=\frac14.
\]
Since
\[
P(A\cap B)=P(A)P(B),
\]
the events are independent.
\end{solution}

\begin{takeaways}
\item Independence is tested by comparing $P(A\cap B)$ with $P(A)P(B)$.
\item Always define the events explicitly before testing.
\item Small sample spaces are often easiest to list completely.
\end{takeaways}

\begin{problem}[Why independence does not depend on order]
In higher studies of probability, statistics, and measure theory, independence is treated as a relationship between events rather than a directional process. So it should not depend on whether we name $A$ first or $B$ first.

Suppose event $B$ is independent of event $A$, meaning that
\[
P(B\cap A)=P(B)P(A).
\]
Prove that event $A$ is independent of event $B$.
\end{problem}

\begin{solution}
Intersection is commutative, so
\[
B\cap A=A\cap B.
\]
Also, multiplication of real numbers is commutative, so
\[
P(B)P(A)=P(A)P(B).
\]

Therefore, from
\[
P(B\cap A)=P(B)P(A),
\]
we immediately get
\[
P(A\cap B)=P(A)P(B).
\]

This is exactly the condition for $A$ to be independent of $B$.

Hence independence is symmetric.
\end{solution}

\begin{takeaways}
\item Independence is a relationship between events, not a directional action.
\item The proof works because both intersection and multiplication are commutative.
\item This symmetry is one reason independence extends naturally to random variables and more advanced probability structures.
\end{takeaways}

\begin{problem}[Expected score of a game]
A spinner gives scores according to the following distribution:
\[
P(X=0)=0.2,\qquad P(X=2)=0.5,\qquad P(X=5)=0.3.
\]
Find the expected score.
\end{problem}

\begin{solution}
The expected value is
\[
E(X)=\sum xP(X=x).
\]
So
\[
E(X)=0(0.2)+2(0.5)+5(0.3)=0+1+1.5=2.5.
\]

Therefore the expected score is
\[
\boxed{2.5}.
\]
\end{solution}

\begin{takeaways}
\item Expected value is a weighted average.
\item Multiply each outcome by its probability, then add.
\item The expectation need not be an actual score on a single play.
\end{takeaways}

\begin{problem}[A fair $n$-sided spinner]
A spinner has $n$ equal sectors labelled
\[
1,2,3,\dots,n.
\]
Let $X$ be the score obtained when the spinner is spun once.

\medskip
\noindent\textit{Example.} For $n=6$, a fair spinner could look like this:

\begin{center}
\begin{tikzpicture}[scale=0.9]
  \def\r{2.1}
  \foreach \i/\lab/\col in {
    0/1/bookpurple!12,
    60/2/yellow!25,
    120/3/bookpurple!20,
    180/4/yellow!35,
    240/5/bookpurple!12,
    300/6/yellow!25
  }{
    \filldraw[fill=\col, draw=black, line width=0.6pt] (0,0) -- (\i:\r) arc (\i:\i+60:\r) -- cycle;
    \node at (\i+30:1.2) {\large $\lab$};
  }
  \draw[line width=0.8pt] (0,0) circle (\r);
  \fill (0,0) circle (0.06);
  \draw[-Latex, line width=1pt, bookpurple] (0,0) -- (78:1.65);
\end{tikzpicture}
\end{center}

\begin{enumerate}[(a)]
  \item Write down $P(X=k)$ for $1\le k\le n$.
  \item Find $E(X)$.
  \item Find $\mathrm{Var}(X)$.
  \item Interpret your result for $E(X)$ in terms of symmetry.
\end{enumerate}

You may use
\[
1+2+\cdots+n=\frac{n(n+1)}{2},
\qquad
1^2+2^2+\cdots+n^2=\frac{n(n+1)(2n+1)}{6}.
\]
\end{problem}

\begin{solution}
\textbf{(a)} Since the spinner is fair, each score is equally likely. Therefore
\[
P(X=k)=\frac1n
\qquad \text{for } k=1,2,\dots,n.
\]

\textbf{(b)} The expected value is
\[
E(X)=\sum xP(X=x)
=\sum_{k=1}^n k\cdot \frac1n
=\frac1n\sum_{k=1}^n k.
\]
Using
\[
\sum_{k=1}^n k=\frac{n(n+1)}{2},
\]
we get
\[
E(X)=\frac1n\cdot \frac{n(n+1)}{2}=\frac{n+1}{2}.
\]

\textbf{(c)} First find
\[
E(X^2)=\sum x^2P(X=x)
=\sum_{k=1}^n k^2\cdot \frac1n
=\frac1n\sum_{k=1}^n k^2.
\]
Using
\[
\sum_{k=1}^n k^2=\frac{n(n+1)(2n+1)}{6},
\]
we have
\[
E(X^2)=\frac{(n+1)(2n+1)}{6}.
\]

So
\[
\mathrm{Var}(X)=E(X^2)-[E(X)]^2
=\frac{(n+1)(2n+1)}{6}-\left(\frac{n+1}{2}\right)^2.
\]
Simplifying,
\[
\mathrm{Var}(X)=\frac{(n+1)(2n+1)}{6}-\frac{(n+1)^2}{4}
=\frac{n^2-1}{12}.
\]

\textbf{(d)} The value
\[
E(X)=\frac{n+1}{2}
\]
is exactly halfway between the smallest and largest possible scores, namely $1$ and $n$. This matches the symmetry of a fair spinner: the distribution is balanced evenly around the centre.
\end{solution}

\begin{takeaways}
\item A fair finite model often leads to a discrete uniform distribution.
\item Uniform distributions are natural examples where symmetry helps predict the mean.
\item Variance can still be found systematically using $E(X^2)-[E(X)]^2$.
\end{takeaways}

\begin{problem}[A random sum from a multiset]
A student has six cards with values
\[
4,\ 5,\ 6,\ 6,\ 7,\ 8.
\]
Three cards are chosen at random without replacement, and their sum is recorded as the random variable $X$.
\begin{enumerate}[(a)]
  \item Construct the probability distribution table for $X$.
  \item Find $E(X)$.
  \item Without calculating $\mathrm{Var}(X)$ fully, explain whether you would expect the spread of this distribution to be larger or smaller than if the six cards were
  \[
  4,\ 5,\ 6,\ 7,\ 8,\ 9.
  \]
  \item Comment briefly on why repeated values in the population affect the shape of the sampling distribution.
\end{enumerate}
\end{problem}

\begin{solution}
\textbf{(a)} Since $3$ cards are chosen from $6$, the number of equally likely selections is
\[
\binom63=20.
\]

Using the repeated value $6$, we count how many selections produce each sum:
\[
\begin{array}{c|c}
\text{Sum} & \text{Selections} \\
\hline
15 & (4,5,6)\text{ in }2\text{ ways} \\
16 & (4,5,7),\ (4,6,6) \\
17 & (4,5,8),\ (4,6,7)\text{ in }2\text{ ways},\ (5,6,6) \\
18 & (4,6,8)\text{ in }2\text{ ways},\ (5,6,7)\text{ in }2\text{ ways} \\
19 & (4,7,8),\ (5,6,8)\text{ in }2\text{ ways},\ (6,6,7) \\
20 & (5,7,8),\ (6,6,8) \\
21 & (6,7,8)\text{ in }2\text{ ways}
\end{array}
\]

So the distribution table is
\[
\begin{array}{c|ccccccc}
x & 15 & 16 & 17 & 18 & 19 & 20 & 21 \\
\hline
P(X=x) & \frac{2}{20} & \frac{2}{20} & \frac{4}{20} & \frac{3}{20} & \frac{4}{20} & \frac{2}{20} & \frac{2}{20}
\end{array}
\]
This simplifies to
\[
\begin{array}{c|ccccccc}
x & 15 & 16 & 17 & 18 & 19 & 20 & 21 \\
\hline
P(X=x) & \frac{1}{10} & \frac{1}{10} & \frac{1}{5} & \frac{3}{20} & \frac{1}{5} & \frac{1}{10} & \frac{1}{10}
\end{array}
\]

\textbf{(b)} The expected value is
\[
E(X)=15\left(\frac{1}{10}\right)+16\left(\frac{1}{10}\right)+17\left(\frac{1}{5}\right)+18\left(\frac{3}{20}\right)+19\left(\frac{1}{5}\right)+20\left(\frac{1}{10}\right)+21\left(\frac{1}{10}\right).
\]
This gives
\[
E(X)=18.
\]

\textbf{(c)} The spread should be \textbf{smaller} than for the cards
\[
4,\ 5,\ 6,\ 7,\ 8,\ 9.
\]
That is because the repeated value $6$ creates more selections with central sums and fewer selections with widely spread values. Repetition pulls the distribution inward.

\textbf{(d)} Repeated values make some totals occur more often than others. This changes the shape of the sampling distribution by increasing clustering around certain sums. In higher studies of probability and statistics, this is an example of how the structure of a population affects the distribution of a statistic computed from a random sample.
\end{solution}

\begin{takeaways}
\item A random variable can be created by applying a rule, here ``add the chosen cards'', to a combinatorial experiment.
\item Equal-likelihood of selections does not mean equal-likelihood of the resulting sums.
\item Repeated values in the underlying population can make the sampling distribution more concentrated around the centre.
\end{takeaways}

\begin{problem}[Selection probability via combinations]
From $6$ girls and $4$ boys, a committee of $3$ students is chosen at random. Find the probability that the committee contains exactly $2$ girls.
\end{problem}

\begin{solution}
The total number of committees is
\[
\binom{10}{3}=120.
\]

To have exactly $2$ girls, choose $2$ of the $6$ girls and $1$ of the $4$ boys:
\[
\binom62\binom41=15\cdot 4=60.
\]

Therefore,
\[
P(\text{exactly 2 girls})=\frac{60}{120}=\frac12.
\]
\end{solution}

\begin{takeaways}
\item When order does not matter, combinations usually give the cleanest probability model.
\item Count favourable selections and total selections in the same style.
\item ``Exactly'' often suggests one precise combination pattern.
\end{takeaways}

\begin{problem}[Leading digit of a height measurement]
An anthropometric study measures the height of randomly selected adults in centimetres.

Let $H$ be the height of a randomly selected adult, and let $X$ be the leading digit of $H$. For example, if
\[
H=175,
\]
then
\[
X=1.
\]

Assume all measured adults are between $100$ cm and $250$ cm tall.
\begin{enumerate}[(a)]
  \item State whether $X$ is a discrete random variable, a continuous random variable, or neither. Justify your answer.
  \item State the realistic sample space for $X$.
  \item A student argues that because almost every adult has a leading height digit of $1$, the value of $X$ is highly predictable and therefore $X$ is not a valid random variable. Evaluate this claim.
\end{enumerate}
\end{problem}

\begin{solution}
\textbf{(a)} The variable $X$ is a \textbf{discrete random variable}.

It is a random variable because its value is determined by a chance process: choosing an adult at random. It is discrete because it can take only separated whole-number values rather than a continuous range.

\textbf{(b)} Since all adults are between $100$ cm and $250$ cm tall, the leading digit can only be
\[
1 \text{ or } 2.
\]
So the sample space for $X$ is
\[
\{1,2\}.
\]

\textbf{(c)} The student's claim is incorrect.

A random variable does \emph{not} need to be uniformly distributed or hard to predict. It only needs to be determined by chance. Even if
\[
P(X=1)
\]
is very large, there is still some chance that
\[
X=2
\]
when a taller adult is selected. Since the exact outcome is not known with certainty before the selection is made, $X$ remains a valid random variable.
\end{solution}

\begin{takeaways}
\item In mathematics, ``random'' means determined by chance, not necessarily evenly spread or highly unpredictable.
\item A discrete random variable may still have a very skewed distribution.
\item This problem also gives a useful contrast with Benford's Law: Benford behaviour often appears in data sets spread across several powers of ten, but adult human heights are trapped in a narrow range, so their leading digits are dominated by $1$.
\end{takeaways}

\begin{problem}[HSC 2023-style: Sample proportion and normal approximation]
A recent census found that $30\%$ of Australians were born overseas. A sample of $900$ randomly selected Australians is surveyed.

Let $\hat p$ be the sample proportion of surveyed people who were born overseas.
\begin{enumerate}[(a)]
  \item Show that
  \[
  \mathrm{Var}(\hat p)=\frac{7}{30000}.
  \]
  \item Use a normal approximation to estimate
  \[
  P(\hat p\le 0.31).
  \]
\end{enumerate}
\end{problem}

\begin{solution}
Let $X$ be the number of surveyed people who were born overseas. Then
\[
X\sim \mathrm{Bin}(900,0.3),
\qquad
\hat p=\frac{X}{900}.
\]

\textbf{(a)} Using the variance formula for a sample proportion,
\[
\mathrm{Var}(\hat p)=\frac{p(1-p)}{n}
=\frac{0.3(0.7)}{900}
=\frac{0.21}{900}
=\frac{21}{90000}
=\frac7{30000}.
\]

\textbf{(b)} Since $n$ is large, we approximate
\[
\hat p \approx N\left(0.3,\frac7{30000}\right).
\]

So
\[
\sigma_{\hat p}=\sqrt{\frac7{30000}}\approx 0.01528.
\]

Now
\[
P(\hat p\le 0.31)\approx P\left(Z\le \frac{0.31-0.30}{0.01528}\right)
=P(Z\le 0.654).
\]

Hence
\[
P(\hat p\le 0.31)\approx 0.74.
\]

So the required probability is approximately
\[
\boxed{0.74}.
\]
\end{solution}

\begin{takeaways}
\item Sample proportion is itself a random variable.
\item Its mean is $p$ and its variance is $\frac{p(1-p)}{n}$.
\item For large samples, normal approximation turns proportion questions into standard normal questions.
\end{takeaways}
